\chapter{Conclusion}
\label{chap:conclusion}

\section{Regime Detection: Status and Interim Conclusion}
\label{sec:conclusion-regime}

This section closes out the regime-detection layer specifically, following the audit, revision, and validation work reported in Sections~\ref{subsec:regime-engine-audit}--\ref{subsec:regime-synthetic-validation} and~\ref{sec:regime-results}. It is scoped deliberately to that layer: the multi-agent deliberation, RL execution, and dashboard layers (Sections~\ref{sec:multi-agent-methodology}--\ref{sec:dashboard-methodology}) were not re-audited as part of this pass, and a full-system conclusion accordingly awaits their own equivalent review.

\subsection{What Was Found}
\label{subsec:conclusion-regime-findings}

Two engineering defects predating this audit were found and corrected: a subsampling procedure in the BIC-based state-count search that fit the model on randomly reordered, non-contiguous rows, corrupting the estimated persistence of every candidate state count; and the absence of any sequence-boundary handling in a pipeline that concatenates multiple assets' histories into one matrix, which had every fit and decode silently treating one asset's final observation as immediately preceding another asset's first. Fixing the second exposed that no seed-sensitivity check had ever existed in the codebase; adding one retroactively showed the production four-asset model's regime partition was not robust to EM initialization (mean pairwise ARI 0.193 across five seeds), a previously unmeasured property of a model that had already been in use.

Independently, a raw-feature variant of the detector was found to have collapsed onto asset identity rather than market regime on the original four-asset panel, traced to a specific, identifiable property of the forex/commodity data source (a literally constant, not merely low-variance, volume-derived feature) dominating an unwhitened Gaussian mixture. This motivated a full panel revision to a crypto-only, five-asset universe, which was verified, not assumed, to resolve the defect (maximum single-asset state share falling from 0.97--0.99 to 0.234) and was carried through the complete engineering and statistical validation battery described above: pipeline validation, empirical state-count re-checking at $n\in\{3,\ldots,8\}$, the six-configuration statistical comparison of Section~\ref{subsec:regime-results-preview}, three independent downstream backtests with an explicit noise control, and a purpose-built synthetic ground-truth validation methodology that, to the best of the author's knowledge, had no precedent elsewhere in this project's prior work.

Two results from that battery are robust across every configuration tested and are the ones this layer's conclusion rests on. The discovered regimes are genuine, statistically significant volatility clusters ($p\approx0$ on realized-volatility distinctness, six for six configurations, no exceptions). And regime conditioning did not demonstrate measurable value in any of the three standalone downstream tests attempted --- volatility-target position sizing, cross-asset systemic allocation, and noise-control-verified directional gating --- nor in next-hour directional prediction, where the regime label alone was consistently at chance level.

\subsection{Interpretation}
\label{subsec:conclusion-regime-interpretation}

These two results are not in tension. A signal can be statistically real without being independently actionable by a simple rule: the synthetic validation work (Section~\ref{subsec:regime-synthetic-validation}) traced the ceiling on per-timestep regime recovery to genuine overlap between the six states' return distributions at hourly resolution --- a Bayes-optimal classifier given the true generating parameters could not exceed an Adjusted Rand Index of 0.047, below what the actual estimated pipeline achieves by aggregating information over rolling windows and exploiting the fitted model's own transition persistence. The states are separable enough for a well-specified model to reproducibly agree on where they are (which the seed-stability results confirm once the feature representation and asset panel were fixed), and volatile enough to be statistically real by any conventional test, without necessarily being separable enough, at single-hour resolution, to hand a simple rule the information it would need to act on them directly.

This bears directly on the core hypothesis motivating this project, stated in Section~\ref{sec:regime-detection}: that explicit regime conditioning improves trading outcomes relative to regime-agnostic learning. The evidence gathered here neither confirms nor refutes that hypothesis at the level it was originally stated --- it refutes a narrower and more easily testable one, that a simple rule conditioned on the regime label outperforms an equivalent rule that is not, under three independent formulations and a genuine noise control. The hypothesis as originally stated is a claim about a \emph{learned} policy's ability to extract value from regime conditioning that a hand-specified rule cannot, which is precisely what the multi-agent and RL layers of Sections~\ref{sec:multi-agent-methodology}--\ref{sec:rl-execution-methodology} were designed to test, and testing it properly requires exercising those layers, not substituting cheaper proxies for them indefinitely. The proxies were worth running first because they are inexpensive and because a negative result at this stage is informative regardless of what the full system later finds: it means the regime vector should be carried forward as one input among several for a learned policy to weigh, not assumed to be a strong prior signal for either that policy or the multi-agent layer's own reasoning.

\subsection{Open Items}
\label{subsec:conclusion-regime-open-items}

Several findings from this pass were deliberately left unresolved rather than silently dropped, and are recorded here for that reason. The production four-asset model's own seed instability was diagnosed but not corrected --- only the revised five-asset panel received the stability improvement, since the four-asset panel was being superseded rather than maintained. The Gaussian emission distribution's underestimation of extreme moves, surfaced independently by both the tail-probability cross-check and by inspection of the fitted variance's departure from an i.i.d.\ approximation (Section~\ref{subsec:regime-synthetic-validation}), was identified but not addressed with an alternative emission distribution, which would be a substantial implementation undertaking whose payoff is uncertain given that the tail-risk gap was diagnosed on the synthetic \emph{calibration} model rather than confirmed on the production fit's own real-data residuals. No component downstream of this layer --- the multi-agent context builder, the RL environment's observation space, or the live dashboard --- was modified to consume the revised five-asset panel; this pass was research-track work in isolated data paths throughout, and a production cutover, if warranted once the downstream layers' own review is complete, remains a separate, explicit decision rather than an implied next step of this chapter.
